\chapter{Introduction au Data Mining}

\section{Pourquoi le Data Mining ? (motivations)}

\begin{Exemple}[le banquier]
Un client demande un crédit. Le banquier veut savoir \emph{s'il va rembourser}. Le client
ne dira jamais \og je ne rembourserai pas \fg{}. Mais le banquier peut \textbf{inférer} le
comportement futur du client à partir de la \textbf{base de données des anciens clients}.
\textbf{C'est exactement ça le Data Mining} : utiliser le passé (les données) pour prédire/
comprendre le futur.
\end{Exemple}

On vit une \textbf{explosion des données} : bases de données, entrepôts, capteurs,
web, e-commerce, hôpitaux, supermarchés… On produit des téraoctets chaque jour, mais
\textbf{seule une infime partie est analysée}. D'où le besoin d'outils capables d'extraire la
\emph{connaissance} cachée dans ces montagnes de données.

\begin{Astuce}
Phrase à retenir du cours : \og \emph{Nous sommes inondés de données, et il nous manque les
connaissances contenues dans ces données.} \fg{}
\end{Astuce}

\section{Qu'est-ce que le Data Mining ?}

\subsection{Trois types de requêtes}
\begin{itemize}[leftmargin=1.4em]
  \item \textbf{Opérationnelle} (base de données / SQL) : \emph{\og Combien de bouteilles de
        vin vendues au 1\textsuperscript{er} trimestre 2010 à la boutique EREVAN ? \fg{}}
  \item \textbf{Analytique} (entrepôt / OLAP) : \emph{\og ventes par département, par marque,
        par trimestre sur 5 ans \fg{}} — agrégation à différents niveaux.
  \item \textbf{Exploratoire} (Data Mining) : \emph{\og Peut-on \textbf{prédire} qu'un client
        va acheter du vin ? Quels produits achète-t-on \textbf{souvent avec} le vin ? \fg{}}
        — impossible en SQL.
\end{itemize}

\begin{Definition}[Data Mining]
\textbf{Ensemble de méthodes et techniques d'analyse de données et d'extraction
d'informations structurées en vue d'aider à la prise de décision.}
Le \og produit \fg{} du Data Mining, ce sont des \textbf{modèles} (les connaissances).
\[ \boxed{\text{Données (information)} \;\xrightarrow{\;\text{Data Mining}\;}\; \text{Modèles (connaissances)}} \]
Historique : \textbf{KDD} (1989, \emph{Knowledge Discovery in Databases}), puis le terme
\textbf{Data Mining} en 1991.
\end{Definition}

\begin{Attention}[: ce que le Data Mining n'est PAS]
Ce n'est ni un entrepôt de données, ni du SQL, ni de l'OLAP, ni de la simple visualisation,
ni de la statistique descriptive, et surtout \textbf{pas une boule de cristal} : il
\emph{extrapole} à partir de données existantes, il ne devine pas magiquement.
\end{Attention}

\section{Les tâches du Data Mining}

C'est \textbf{le tableau le plus important du chapitre}. Tout le reste du cours détaille ces
tâches.

\begin{center}
\renewcommand{\arraystretch}{1.35}
\begin{tabular}{>{\bfseries}l p{6.2cm} l c}
\toprule
Tâche & Objectif & Méthodes & Apprentissage\\
\midrule
Classification & Attribuer une \emph{classe connue} à un objet (crédit oui/non) &
  Arbres, K-NN, réseaux de neurones, Bayes & supervisé\\
Prédiction & Estimer une \emph{valeur future} d'un attribut &
  Arbres, réseaux de neurones & supervisé\\
Association & Trouver les attributs \emph{corrélés} (panier : poisson \& vin) &
  Apriori, FP-Growth & non supervisé\\
Segmentation & Former des \emph{groupes homogènes} (classes inconnues) &
  K-means, K-NN, réseaux de neurones & non supervisé\\
\bottomrule
\end{tabular}
\end{center}

\begin{Definition}[Supervisé vs non supervisé]
\textbf{Supervisé} : l'apprenant reçoit des exemples avec \emph{entrée ET sortie} (la classe
est connue) $\rightarrow$ classification, prédiction.\\
\textbf{Non supervisé} : seulement des \emph{entrées}, pas de classe $\rightarrow$ association,
segmentation.
\end{Definition}

On distingue aussi les \textbf{modèles prédictifs} (extrapolent une info nouvelle :
classification, prédiction) et les \textbf{modèles descriptifs} (mettent en évidence une info
noyée : segmentation, association).

\section{Le processus CRISP-DM}

\textbf{CRISP-DM} (\emph{Cross Industry Standard Process for Data Mining}) est la
méthodologie la plus utilisée. Elle découpe un projet en \textbf{6 phases}, dans un cycle
où l'on fait des \og aller-retour \fg{}.

\begin{center}
\renewcommand{\arraystretch}{1.25}
\begin{tabular}{c >{\bfseries}l p{9.5cm}}
\toprule
\# & Phase & En quoi ça consiste\\
\midrule
1 & Compréhension du métier & Objectifs de l'entreprise $\rightarrow$ problème de data mining.\\
2 & Compréhension des données & Recueillir, explorer, évaluer la qualité.\\
3 & Préparation des données & Nettoyer, sélectionner, transformer (\textbf{$\approx$60\,\% du temps}).\\
4 & Modélisation & Choisir les algorithmes et calibrer leurs paramètres.\\
5 & Évaluation & Les résultats répondent-ils aux objectifs ?\\
6 & Déploiement & Mettre en œuvre les décisions (côté maître d'ouvrage).\\
\bottomrule
\end{tabular}
\end{center}

\begin{Exemple}[étude de cas Daimler-Chrysler]
Objectif : \textbf{analyser les plaintes clients} dans l'automobile (réduire les coûts,
améliorer la satisfaction). Les chercheurs ont suivi CRISP-DM : compréhension métier
(\og y a-t-il un lien entre plaintes ? \fg{}), exploration du système QUIS (7~millions de
véhicules, 40~Go), préparation \textbf{très longue} (formats incompatibles), modélisation
(réseaux bayésiens + règles d'association), évaluation \textbf{décevante} (modèles peu
fiables à cause de la structure de la base), déploiement limité à un projet pilote. \textbf{Leçon :
ne jamais sous-estimer la préparation des données.}
\end{Exemple}

\section{Outils et sources de données}
Logiciels phares (sondages KDnuggets) : \textbf{RapidMiner}, \textbf{R}, \textbf{Excel},
\textbf{Weka}, \textbf{Python} (numpy/pandas/scikit-learn), SAS, KNIME… Dépôts de données
pour s'entraîner : \textbf{UCI} (\texttt{archive.ics.uci.edu}), \textbf{Kaggle},
\textbf{Hugging Face}.

\begin{Resume}
Le Data Mining = extraire des \textbf{modèles} décisionnels des données. 4 tâches
(classification, prédiction, association, segmentation), 2 régimes (supervisé / non
supervisé), 1 méthodologie reine (\textbf{CRISP-DM}, 6 phases) où la \textbf{préparation des
données} domine. Les chapitres suivants détaillent : la préparation (chap.~2), la
classification (chap.~4) et la segmentation (chap.~5).
\end{Resume}
